\documentclass[conference]{IEEEtran}
\usepackage[utf8]{inputenc}
\usepackage[T1]{fontenc}
\usepackage{cite}
\usepackage{amsmath,amssymb}
\usepackage{graphicx}
\usepackage{booktabs}
\usepackage{url}

\title{Effectiveness of a Generate--Verify--Repair Pipeline with Batfish for LLM‑Generated Vendor CLI Configurations}

\author{
\IEEEauthorblockN{Autonomous AI Researcher}
\IEEEauthorblockA{CNSM 2026 Agentic AI Researcher Track}
}

\begin{document}

\maketitle

\begin{abstract}
We report a preregistered paired experiment (cnsms2026-001-gvr-batfish-prereg-v1) that evaluated whether a generate--verify--repair (GVR) pipeline---shared initial LLM candidate (gpt-5-mini) + a deterministic Batfish‑style validator + at most one automated LLM repair---reduces deterministic‑verifier detected semantic/safety failures on intent\ensuremath{\rightarrow}vendor‑CLI tasks. Planned N=300 pairs; 266 complete pairs were analysed after preregistered exclusions. Baseline pass-rate = 260/266 (97.744\%); guarded (final GVR) pass-rate = 265/266 (99.624\%). Paired difference = 0.01879699248; exact McNemar two‑sided p = 0.0625. Paired bootstrap (B=10,000, seed=1729) 95\% percentile CI = [0.0037593984962406013, 0.03759398496240601]. Repair was invoked for 6 tasks; median repair iterations per task = 0. Per‑episode latency was not recorded in per‑task artifacts and thus could not be evaluated. Under shared\_initial\_candidate semantics the observed absolute improvement is small and did not meet the preregistered \ensuremath{\alpha}=0.05 decision rule. We provide preregistered design details, deterministic validator semantics, exact paired counts, representative artifact‑grounded repair examples, contamination findings (no exact public duplicates flagged), and a focused discussion of limitations and implications.
\end{abstract}

\section{Introduction}
Generative large language models (LLMs) are actively explored for NetOps tasks such as intent\ensuremath{\rightarrow}CLI translation, zero‑touch configuration synthesis, and automated maintenance workflows [1], and surveys emphasize both opportunity and risk when LLMs are applied to operational network configuration [2]. Stochastic LLM outputs can produce syntactic or semantic violations---missing required commands, unsafe command orderings, or constraint breaches---that create operational risk if applied directly. To mitigate that risk, pipelines that interpose deterministic verifiers (e.g., Batfish or header‑space analyses) and bounded repair loops---collectively generate--verify--repair (GVR)---have been proposed to provide deterministic acceptance criteria and automated correction [3,4]. Prior work presents components and prototypes, but preregistered, paired evaluations that archive verifier traces and quantify verifier‑triggered repair benefit on reproducible NetOps task banks remain limited.

This study tested the preregistered question (cnsms2026-001-gvr-batfish-prereg-v1): Does a GVR pipeline that uses a single sampled initial LLM candidate per task (shared\_initial\_candidate semantics), a deterministic Batfish‑style validator, and at most one automated LLM repair call reduce deterministic‑verifier detected semantic/safety failures relative to the same initial candidate without repair? Our principal contribution is a preregistered, artifact‑archived paired evaluation executed under the CNSM Agentic AI Researcher constraints that reports exact paired counts, inferential outcomes, execution accounting, representative artifact examples, and an evidence‑grounded interpretation.

\section{Related Work}
Prior work relevant to this study spans (a) applied LLM prototypes and benchmarks for network configuration, (b) surveys documenting LLM failure modes and recommendations to ground generation, (c) proposals for generate--verify--repair and related architectures for deterministic validation, and (d) established deterministic network verification tools that serve as practical oracles in automated pipelines. We synthesize those threads and emphasize aspects most relevant to a preregistered paired GVR evaluation.

LLM-driven NetOps prototypes and benchmarks: Recent system‑level and benchmark efforts investigate whether LLMs can aid network configuration generation and automation. NetConfEval and similar task collections examine intent\ensuremath{\rightarrow}CLI translation and measure correctness under vendor‑dialect constraints; these works provide public task exemplars and motivate reproducible task generators for controlled evaluation [1]. Such benchmark tasks typically encode required command sequences and vendor syntax, which directly map to the validator check types we archive in this study (required‑sequence, allowed‑field checks, topology reachability). Importantly, prior benchmarks often evaluate single‑pass generation accuracy rather than integrated verifier‑triggered repair, which motivates the paired GVR design adopted here.

Failure modes, grounding, and mitigation recommendations: Survey and synthesis papers summarize that LLMs commonly produce both syntactic and semantic errors when applied to domain‑constrained tasks: missing commands, incorrect ordering, and hallucinated or unsafe modifications [2]. Those surveys recommend grounding strategies (retrieval, tool use), constrained generation, and verifier insertion to reduce the operational risk of accepting model outputs without deterministic checks. Our work adopts those prescriptions by interposing a deterministic validator and bounding the repair loop to one automated correction attempt, aligning experimentally with the practical mitigations advocated in the literature.

Generate--verify--repair and bounded repair proposals: The idea of interposing deterministic oracles and using closed‑loop repair is well established in proposals for regulated or safety‑sensitive domains. Convergent GVR architectures aim to combine stochastic generation with deterministic acceptance criteria and limited repair to achieve practical correctness guarantees while preserving the generative model's productivity [3]. These proposals emphasize publishing verifier traces and repair prompts so that corrected outputs are auditable---an archival design we implemented in the execution artifacts (validator traces, repair\_prompt\_sha256, repair provider traces).

Deterministic network verifiers as oracles and their coverage limits: Tools like Batfish and header‑space analyses provide deterministic, rule‑based checks for reachability, ordering, and configuration invariants and have been incorporated into research prototypes that integrate model outputs with formal validation [4]. Batfish‑style validators excel at expressing and mechanically checking reachability and required‑sequence constraints on modeled topologies; however, their guarantees are strictly limited to the rule battery they encode. This distinction matters experimentally: a verifier that checks an explicit finite set of syntactic/semantic invariants yields a binary pass/fail endpoint that is reproducible and auditable, but it cannot detect semantic failures outside its coverage. Our deterministic\_netops\_validator\_v5 follows this same engineering trade‑off (syntactic parsing, required‑sequence checks, allowed‑fields enforcement, and topology reachability), and we archive per‑episode validation\_before/after traces so readers can inspect exactly which rules were exercised on each episode.

What LLMs need to synthesize correct router configurations and constrained decoding: Empirical studies analyzing LLM difficulty with router configuration point to sequencing, vendor‑specific syntax, and dependency constraints as recurring challenges [5]. These analyses motivate two complementary mitigations: (1) stronger first‑pass constraints or constrained decoding to reduce syntactic/format errors, and (2) verifier‑triggered repair to correct semantic violations that slip through first pass. Our study deliberately fixes the initial candidate (shared\_initial\_candidate) and thereby isolates the marginal effect of the verifier‑triggered repair step rather than conflating improvements due to constrained first‑pass prompting or multiple draws.

Gaps in prior evaluations that motivated this study: The supplied literature shows active proposals and component evaluations but comparatively fewer preregistered, paired experiments that (i) use shared\_initial\_candidate semantics to measure the isolated effect of bounded repair, (ii) archive raw model traces together with deterministic validator traces and repair provider traces for each episode, and (iii) apply explicit exclusion, retry, and contamination rules in a confirmatory design. Those gaps motivated our preregistered paired design and the artifact‑heavy archival choices (manifested in the execution and analysis artifacts we publish alongside this manuscript).

Methodologically relevant syntheses for experimental design: From the surveyed and prototyped prior work we derived three concrete design choices: (1) use a deterministic validator battery (Batfish‑style) to obtain an auditable binary endpoint; (2) bound the repair loop (<=1 automated repair call per failed candidate) to maintain a clear paired estimand under shared\_initial\_candidate semantics; and (3) preregister explicit missingness, retry, and contamination treatments to avoid post hoc analytic choices. Each of these choices has precedent or justification in the cited literature: verifiers as deterministic oracles [4], GVR architectures for regulated settings [3], and benchmark/task construction recommendations that emphasize reproducibility and contamination controls [1,2].

In sum, our related‑work synthesis locates this study at the intersection of applied intent\ensuremath{\rightarrow}CLI benchmarks, survey‑driven mitigation recommendations, GVR methodological proposals, and deterministic verification tooling. The specific empirical contribution is a preregistered, archived paired GVR evaluation under shared\_initial\_candidate semantics that fills a documented gap in artifact‑grounded confirmatory assessments of verifier‑triggered repair effectiveness.

\section{Methodology}
Preregistration and estimand: The study preregistration is cnsms2026-001-gvr-batfish-prereg-v1 (manifest SHA‑256: 46a22349\allowbreak{}2f760950\allowbreak{}b3d06f47\allowbreak{}5129ce21\allowbreak{}e52e296c\allowbreak{}4c2a1e09\allowbreak{}df3ccd50\allowbreak{}bec9f891). Primary estimand: paired\_success\_rate\_difference\_guarded\_minus\_baseline under shared\_initial\_candidate semantics. Under those semantics the identical single sampled initial model candidate is used to compute both outcomes: baseline single‑shot candidate \ensuremath{\rightarrow} deterministic validator pass/fail, and guarded outcome \ensuremath{\rightarrow} same candidate; if it fails, the guarded arm may invoke at most one automated LLM repair call and the final validator pass/fail is recorded. This isolates the effect of the allowed validator‑triggered repair step for a fixed initial candidate.

Task bank and sampling: The preregistered task bank deterministically produced 300 intent\ensuremath{\rightarrow}CLI pairs (150 Cisco‑like; 150 Juniper‑like) using a seeded synthetic task generator combined with public NetConfEval‑style examples; per‑task generator ids, seeds, and manifest entries are archived. Inclusion/exclusion, retry, and missingness rules followed the preregistration: provider/API transient failures allowed up to two automatic retries; pairs where both arms failed due to provider/infrastructure errors were excluded from the primary analysis. The execution manifest records planned\_episode\_count = 600 and completed\_episode\_count = 532; missingness and reconciliation artifacts are archived.

Model, orchestration, and sampling disclosure: The declared model was gpt-5-mini (execution/model\_configuration.json field declared\_model\_version = "gpt-5-mini"). Archived configuration fields include max\_output\_tokens = 1024, maximum\_attempts\_per\_call = 1, provider = openai\_responses, and temperature = null (unset). Provider accounting (deterministic\_reconciliation.json) records hosted\_model\_call\_event\_count = 306, completed\_hosted\_model\_call\_event\_count = 272, and failed\_hosted\_model\_call\_event\_count = 34. Because temperature was unset and no centralized deterministic sampling seed was archived for provider calls, nondeterministic sampling cannot be excluded for recorded model calls; we note this explicitly in the Disclosure and Limitations.

Deterministic validator semantics and scoring rules: The binary oracle used across episodes was deterministic\_netops\_validator\_v5. For each candidate the validator performed: (1) CLI syntactic parsing/normalization; (2) per‑task required‑sequence checks (expected\_sequence vs parsed commands); (3) constrained‑field touch checks (allowed\_fields, allowed\_touched\_fields); and (4) reachability/constraint validation on the synthetic topology model. Each episode archives validation\_before and validation\_after traces containing parsed\_command\_count, required\_command\_count, observed\_touched\_field\_count, violation codes, and normalized\_configuration; the validator's pass/fail output is the primary binary endpoint. The confirmatory scoring rule was 'full\_intent\_constraint\_validation' as recorded in analysis artifacts.

Repair loop mechanics (bounded, archived): Per the preregistration the repair stage is bounded to at most one automated LLM repair call per task when the initial candidate fails the validator. The repair prompt includes validator failure codes and task constraints and requests a corrected configuration; repair provider traces and repair\_prompt\_sha256 are archived when invoked. Stage accounting shows shared\_initial\_generation = 300 and repair = 6 (six repair calls executed). Repair artifacts and validator traces for repaired episodes are included in execution/scoring and representative artifacts described below.

Analysis plan and inferential methods: The prespecified primary test is the exact McNemar two‑sided test on paired binary outcomes at \ensuremath{\alpha}=0.05. Interval estimation was preregistered as a paired bootstrap percentile CI (B=10,000, seed=1729) for the paired difference to quantify magnitude and uncertainty. Missingness handling followed the preregistered 'complete\_pair\_primary' rule: both‑arm provider failures were excluded from the primary test; archived sensitivity analyses treat excluded episodes conservatively as failures. Secondary estimands included median\_repair\_iterations\_per\_task and median end‑to‑end latency; per‑episode latency was not archived and thus could not be evaluated. We included the paired bootstrap CI alongside McNemar because exact McNemar p‑values can be discrete and conservative with few discordant pairs, whereas a bootstrap CI summarizes empirical paired‑difference variability (see Results for the observed small‑discordant behavior).

\section{Results}
Execution accounting and sample construction. The preregistered run planned 600 model‑condition episodes (300 tasks \ensuremath{\times} 2 conditions). The orchestrated pipeline completed 532 episodes and recorded 68 failed episodes (execution\_manifest.json). After applying preregistered exclusion rules that remove pairs affected by both‑arm provider/infrastructure failures, the confirmatory analysis used 266 complete paired units (analysis/missingness\_summary.csv). Key provider‑call accounting from the deterministic reconciliation artifact: hosted\_model\_call\_event\_count = 306 total model call events, of which 272 completed successfully and 34 failed; cache was unused (cache\_hit\_count = 0, cache\_miss\_count = 272). Stage accounting shows that a shared initial generation was produced for every task (shared\_initial\_generation = 300) and the automated repair stage executed in 6 episodes (repair = 6). All execution‑level manifests and per‑episode traces are archived to permit replay and audit.

Primary paired outcomes and exact counts. The confirmatory paired contingency (guarded vs baseline) computed on the 266 complete pairs is shown in analysis/paired\_contingency\_table.csv and summarized here: n11 = 260 (both baseline and guarded pass), n10 = 5 (guarded pass, baseline fail), n01 = 0 (baseline pass, guarded fail), n00 = 1 (both fail). These margins yield observed success rates baseline = 260/266 = 0.9774436090 and guarded = 265/266 = 0.9962406015 with an absolute paired difference of 0.01879699248.

Inferential results: exact McNemar and paired bootstrap CI. Per the preregistered primary test, the exact McNemar two‑sided p‑value for the observed discordant counts (n10=5, n01=0) is p = 0.0625, which does not meet the preregistered \ensuremath{\alpha} = 0.05 decision rule for rejection. To quantify effect magnitude uncertainty we computed the preregistered paired bootstrap percentile 95\% CI with B = 10,000 resamples and seed = 1729 (analysis/analysis\_log.jsonl); the resulting CI is [0.0037593984962406013, 0.03759398496240601], narrowly excluding zero. Reporting both results preserves the preregistered decision rule while also conveying interval evidence: the discrete nature of McNemar's exact two‑sided test with few discordants can produce conservative p‑values even when bootstrap resampling indicates a small positive paired difference with limited sampling variability. Both inferential artifacts (exact p and bootstrap CI with documented seed and resamples) are archived.

Repair invocation behavior and representative repaired episodes. The repair stage executed for 6 of the 300 tasks (stage\_counts: repair = 6). Median repair iterations per task across the analysed sample is 0 (majority of tasks required no repair because the shared initial candidate already passed validation). Representative repaired episode: task-000008. The shared initial candidate for task-000008 (execution/responses/task-000008-shared-initial.txt) contained a truncated final token and failed validator checks (validation\_before showed parsed\_command\_count < required\_command\_count). The automated repair call produced a corrected candidate (execution/responses/task-000008-guarded.txt) whose validator trace (execution/scoring/task-000008-guarded.json) records validation\_after.valid = true with no violation codes. Per‑episode scoring JSONs record whether a repair was applied, the repair\_prompt reference, and the final normalized\_configuration accepted by the deterministic validator; all such per‑episode repair traces are archived for examination.

Contamination and sensitivity to missingness. The preregistered contamination detector flagged no exact public duplicates (analysis/contamination\_summary.csv empty), so no tasks were excluded from the primary analysis on contamination grounds. The preregistered missingness policy excluded 34 both‑arm failed episodes from the primary test; a sensitivity analysis that conservatively treats those 34 excluded pairs as failures (i.e., complete N = 300 with both arms set to failure where missing) yields baseline = 260/300 = 0.866667 and guarded = 265/300 = 0.883333, but retains the same discordant structure (n10=5, n01=0) leading to the same exact McNemar p = 0.0625. Thus, under the preregistered conservative missingness treatment the formal decision does not change; archived missingness\_summary.csv and deterministic\_reconciliation.json provide the complete accounting.

Supplementary quantitative artifacts and reproducibility. Analysis artifacts included in the evidence bundle and referenced above provide the reproducible numeric foundations for these results: analysis/condition\_summary.csv (baseline: 260 successes, 6 failures; guarded: 265 successes, 1 failure), analysis/paired\_contingency\_table.csv (n11, n10, n01, n00 table), analysis/missingness\_summary.csv (counts of excluded and failed episodes), and analysis/analysis\_log.jsonl (bootstrap settings: B=10,000, seed=1729). deterministic\_reconciliation.json documents provider call totals and stage counts that explain model\_calls\_used = 306 and repair\_calls = 6. All files and their hash manifests are archived to enable exact reanalysis or targeted follow‑up audits.

\section{Discussion}
Statistical interpretation and small‑discordant behavior: The guarded GVR pipeline produced a small absolute improvement (\textasciitilde{}1.88 percentage points). Under the preregistered exact McNemar decision rule the improvement did not reach \ensuremath{\alpha}=0.05 (p=0.0625). The paired bootstrap CI that excludes zero highlights finite‑sample discreteness: with few discordants (n10=5, n01=0) McNemar's discrete two‑sided p‑value can be conservative for small effects while bootstrap percentile CIs reflect resampled paired‑difference variability. We report both inferential perspectives so readers may weigh formal hypothesis rejection against effect magnitude and CI width.

Operational implications and boundary conditions grounded in artifacts: Baseline first‑pass accuracy was high (\ensuremath{\approx}97.7\%), leaving few failures for the repair stage to correct---hence rare repair invocations (6/300). Stage accounting and provider traces show repair-stage overhead was small in the observed run, but we did not archive per‑episode latency and therefore cannot quantify repair latency or end‑to‑end timing against the preregistered operational bound (median <= 60 s). The low repair rate suggests bounded automated repair adds limited marginal benefit on high‑first‑pass corpora; however, GVR may yield larger absolute improvements when first‑pass accuracy is lower, verifier coverage expands, or when the verifier detects failure modes common in production workloads.

Threats to validity (artifact‑grounded): Internal validity: the shared\_initial\_candidate estimand measures the effect of validator‑triggered repair for a fixed initial sample and does not assess independent multiple draws or constrained first‑pass prompting because independent first‑pass arms were not executed under the preregistration. Construct validity: deterministic\_netops\_validator\_v5 enforces a finite, archived rule battery (syntactic parsing, required\_sequence, constrained\_field checks, reachability) and therefore semantic violations outside that battery are undetected and unmeasured. External validity: tasks derive from public NetConfEval‑style examples and a seeded synthetic generator limited to two vendor dialects; results may not generalize to other vendors, larger topologies, or production heterogeneous workloads. Conclusion validity: the loss of 34 both‑arm provider failures reduced effective N and the observed few discordant pairs limit power for small effect detection.

Design lessons and recommendations grounded in execution artifacts: (1) When baseline pass rates are high, design experiments targeting lower‑baseline corpora or include arms that modify first‑pass sampling (e.g., multiple draws or constrained prompts) to increase discordant counts and power. (2) Archive per‑episode pipeline timing (request\ensuremath{\rightarrow}generation\ensuremath{\rightarrow}validation\ensuremath{\rightarrow}repair\ensuremath{\rightarrow}final validation) to assess operational latency bounds; our run lacked per‑episode latency and thus the preregistered latency estimand could not be evaluated. (3) Expand and publish validator rule coverage and include tests that exercise verifier false negatives; per‑episode validator traces permit such audits and should be part of future preregistrations. (4) Consider paired designs that explicitly compare shared\_initial\_candidate semantics to independent generation semantics in separate preregistered arms; such contrasts require independent initial draws and were outside this run by design.

Reproducibility and archival resources: All raw responses, execution/provider call traces, validator traces, scoring JSONs, analysis artifacts (paired contingency table, condition summary, missingness summary, contamination summary), deterministic\_reconciliation.json, and model\_configuration.json are archived and hashed in the evidence bundle. The archived analysis\_log.jsonl records bootstrap settings (B=10,000, seed=1729). These artifacts permit exact reexecution of the confirmatory analysis and targeted replay of repaired episodes.

\section{Limitations}
\begin{itemize}
\item Shared\_initial\_candidate semantics: both arms used the same single sampled initial candidate per preregistration; the estimand measures validator‑triggered repair benefit for that fixed initial sample and does not evaluate independent first‑pass generations or constrained first‑pass prompting strategies.
\item Missingness and sample reduction: planned N=300 pairs; after infrastructure/provider exclusions 266 pairs were complete and 34 both‑arm episodes were excluded per preregistered rules, reducing effective sample size and precision.
\item Small discordant counts: primary inference used exact McNemar with few discordants (n10=5, n01=0); conclusions are sensitive to inferential choice and limited power.
\item Validator coverage: deterministic\_netops\_validator\_v5 enforces a finite set of syntactic/semantic/reachability rules; semantic violations outside its coverage are possible and unmeasured.
\item Runtime nondeterminism and sampling: archived execution/model\_configuration.json records temperature = null and no deterministic seed; nondeterministic sampling cannot be excluded for recorded model calls.
\item H2 latency not evaluated: per‑episode end‑to‑end latency was not present in archived per‑task artifacts and therefore the preregistered latency bound (median <= 60 s) could not be assessed.
\item External validity: task corpus derives from public NetConfEval‑style examples and a seeded synthetic generator limited to two vendor‑dialect clusters; results do not imply universal benefit across other vendors, larger topologies, or production deployments.
\end{itemize}

\section{Conclusion}
We executed a preregistered paired experiment that evaluated a GVR pipeline (gpt‑5‑mini + deterministic\_netops\_validator\_v5 + <=1 automated repair) on intent\ensuremath{\rightarrow}CLI tasks under shared\_initial\_candidate semantics. Across 266 analysed pairs the guarded pipeline improved validator pass rate by 0.01879699248 but did not meet the preregistered exact‑McNemar \ensuremath{\alpha}=0.05 decision rule (p=0.0625); a paired bootstrap CI (B=10,000, seed=1729) narrowly excludes zero. Repair calls were rare (6/300) and median repair iterations per task = 0. Per‑episode latency was not archived and could not be assessed. All execution and analysis artifacts (raw responses, validator traces, provider call accounting, and reconciliation manifests) are archived in the evidence bundle to support reanalysis and follow‑on work.

References
[1] C. Wang, M. Scazzariello, A. Farshin, S. Ferlin, D. Kostić, and M. Chiesa, "NetConfEval: Can LLMs Facilitate Network Configuration?," in Proc. ACM, 2024, doi: 10.1145/3656296.
[2] S. Y. Long, J. Tan, B. Mao, F. Tang, Y. Li, M. Zhao, and N. Kato, "A Survey on Intelligent Network Operations and Performance Optimization Based on Large Language Models," IEEE Commun. Surveys \& Tutorials, 2025, doi: 10.1109/comst.2025.3526606.
[3] R. Cadenas, "Convergent Correctness in Stochastic Code Generation: A Generate‑Verify‑Repair Architecture for Deterministic Validation of Autonomous AI Coding Agents in Regulated Environments," SSRN, 2026, doi: 10.2139/ssrn.6754899.
[4] M. Brown, A. Fogel, D. Halperin, V. Heorhiadi, R. Mahajan, and T. Millstein, "Lessons from the evolution of the Batfish configuration analysis tool," Proc., 2023, doi: 10.1145/3603269.3604866.
[5] R. Mondal, A. Tang, R. Beckett, T. Millstein, and G. Varghese, "What do LLMs need to Synthesize Correct Router Configurations?," Proc., 2023, doi: 10.1145/3626111.3628194.

\section*{Disclosure Statement}
Disclosure (concise): Declared model: gpt-5-mini (execution/model\_configuration.json archived; declared\_model\_version = "gpt-5-mini"). Orchestration adapter: hosted\_netops\_gvr\_v1. Deterministic validator: deterministic\_netops\_validator\_v5. Preregistration: cnsms2026-001-gvr-batfish-prereg-v1 (manifest SHA‑256: 46a22349\allowbreak{}2f760950\allowbreak{}b3d06f47\allowbreak{}5129ce21\allowbreak{}e52e296c\allowbreak{}4c2a1e09\allowbreak{}df3ccd50\allowbreak{}bec9f891). Initial master prompt immutable reference: provenance/master\_prompt.txt (SHA‑256: 1872df1e\allowbreak{}1805d2d9\allowbreak{}6940456c\allowbreak{}a016bd66\allowbreak{}5d1d5196\allowbreak{}add77f5a\allowbreak{}cdf1582b\allowbreak{}b39b15ba). Human role and autonomy boundary: before the autonomy lock a human Research Director selected the topic family and approved pre‑lock infrastructure development; after the autonomy lock the registered pipeline autonomously performed literature selection, preregistration, experimental execution, analysis, manuscript drafting, AI peer review simulation, and revision. Nondeterminism disclosure: archived execution/model\_configuration.json records temperature = null (unset) and no deterministic sampling seed is archived; therefore nondeterministic sampling of model outputs cannot be excluded for the recorded provider calls. Execution and analysis artifacts, logs, and hash manifests are archived in the evidence bundle referenced in the preregistration.

\begin{thebibliography}{99}
\bibitem{ref1} Changjie Wang, Mariano Scazzariello, Alireza Farshin, Simone Ferlin, Dejan Kostić, Marco Chiesa, "NetConfEval: Can LLMs Facilitate Network Configuration?", 2024, 10.1145/3656296
\bibitem{ref2} S.Y. Long, Jingjing Tan, Bomin Mao, Fengxiao Tang, Yangfan Li, Ming Zhao, Nei Kato, "A Survey on Intelligent Network Operations and Performance Optimization Based on Large Language Models", 2025, 10.1109/comst.2025.3526606
\bibitem{ref3} Rafael Cadenas, "Convergent Correctness in Stochastic Code Generation: A Generate-Verify-Repair Architecture for Deterministic Validation of Autonomous AI Coding Agents in Regulated Environments", 2026, 10.2139/ssrn.6754899
\bibitem{ref4} Matt Brown, Ari Fogel, Daniel Halperin, Victor Heorhiadi, Ratul Mahajan, Todd Millstein, "Lessons from the evolution of the Batfish configuration analysis tool", 2023, 10.1145/3603269.3604866
\bibitem{ref5} Rajdeep Mondal, Alan Tang, Ryan Beckett, Todd Millstein, George Varghese, "What do LLMs need to Synthesize Correct Router Configurations?", 2023, 10.1145/3626111.3628194
\end{thebibliography}

\end{document}
